\begin{table}[htbp]
\centering
\begin{threeparttable}
\caption{Run\_10 execution-quality diagnostics and model robustness indicators.}
\label{tab:qualitydiag}
\small
\setlength{\tabcolsep}{4.5pt}
\renewcommand{\arraystretch}{1.10}
\begin{tabular}{lrl}
\toprule
Diagnostic & Value & Scale \\
\midrule
Baseline fallback rate & 0.000 & \% \\
Post fallback rate & 0.000 & \% \\
Attack heuristic pass rate & 0.000 & \% \\
Post heuristic pass rate & 0.000 & \% \\
ICC(1) |Delta| & 0.883 & ratio \\
Ridge CV-R\textasciicircum 2 & -0.134 & R\textasciicircum 2 \\
Random forest OOB R\textasciicircum 2 & -0.110 & R\textasciicircum 2 \\
Elastic-net selected features & 0.000 & count \\
\bottomrule
\end{tabular}
\begin{tablenotes}[flushleft]
\footnotesize
\item Note. These diagnostics are intended to prevent overinterpretation. High fallback rates or near-zero predictive fit imply that downstream susceptibility coefficients are methodological diagnostics rather than substantive evidence about human-like persuasion dynamics.
\end{tablenotes}
\end{threeparttable}
\end{table}